\section{独立監査 C：Class 5 古典 $r$-matrix / monodromy route}
\subsection{目的と今回の結論}

2026年のClass 5原論文は、量子Lax作用素から古典空間Lax行列を構成し、Kostant--Kirillov型のPoisson bracketと古典$r$-matrixを得た一方、運動方程式を再現するcompanion matrixを見つけられなかったと記している\cite{DFN}。原論文はその代わりにboost法から保存量のtowerを構成して古典可積分性を示した。

先行する我々の計算報告では、原論文とは別の方向から、有限格子の量子零曲率方程式に現れる時間Lax作用素を直接連続極限へ移した。その際、同じ物理格子点を共有する作用素積のlower symbolを因子化すると有限の項を失うことを確認し、その補正を保持することで時間Lax行列を得た。さらに零曲率方程式とHamiltonianからの運動方程式がoff-shellで同値であることを確認した。

残っていた確認は、その時間Lax行列が、古典可積分系で標準的な$r$-matrix／monodromy routeから得る時間成分と一致するかどうかであった。Doikou--Karaiskosの記法では、transfer matrixを$t(\zeta)$、空間Lax行列を$\mathbb U(x,\mu)$とすると、生成Hamiltonian flowは
\begin{equation}
\{\ln t(\zeta),\mathbb U(x,\mu)\}
=
\partial_x\mathbb V(x;\zeta,\mu)
+[\mathbb V(x;\zeta,\mu),\mathbb U(x,\mu)]
\label{audit5r:eq:STSflow}
\end{equation}
と書け、時間側の生成行列$\mathbb V$はmonodromyと古典$r$-matrixから構成される\cite{DK}。この種の関係を以下ではSemenov--Tian--Shansky型の構成と呼ぶ。

今回の結論は、Class 5でこの照合が通った、というものである。二つの時間行列は場に依存しないスカラー行列を除いて一致する。従って、Class 5/6の量子--古典Lax bridgeについて、追加の模型計算を必須条件として残す必要はない。

\subsection{記号と物理的意味}

本報告では、monodromyを展開するためのスペクトルパラメータと、最終的なLax pairのスペクトルパラメータを区別する。この区別をしないと、式の意味が分かりにくくなるためである。

\begin{longtable}{L{39mm}L{105mm}}
\toprule
記号 & 物理的意味と定義\\
\midrule
\endhead
古典spin場 $\Svec(x)$ & 単位球面上の古典spin密度。成分を$S^1,S^2,S^3$とする。\\
複素spin成分 $p,q,z$ & $p:=S^1+\ii S^2$, $q:=S^1-\ii S^2$, $z:=S^3$。従って$\Svec^2=1$は$pq+z^2=1$となる。$p$は運動量ではない。\\
Laxスペクトルパラメータ $\mu$ & 最終的な空間Lax行列$U(x,\mu)$と時間Lax行列$V(x,\mu)$に残るスペクトルパラメータ。\\
生成スペクトルパラメータ $\zeta$ & transfer matrix $t(\zeta)$を展開し、可積分hierarchyの各Hamiltonian flowを取り出すための別のスペクトルパラメータ。\\
coherent-state projector $\rho$ & spin-$\tfrac12$ coherent stateに対応するrank-one行列。式\eqref{audit5r:eq:rho}で定義する。\\
Jordanian有限項 $K$ & Class 5空間Lax行列のうち、$\mu\to0$で$1/\mu$極を持たない変形項。式\eqref{audit5r:eq:UK}で定義する。\\
permutation作用素 $P_{12}$ & 二つの$2$次元補助空間を交換する$4\times4$行列。\\
上向き基底へのprojector $\Pup$ & $\Pup:=\operatorname{diag}(1,0)$。\\
raising matrix $\Eplus$ & $\Eplus:=\begin{psmallmatrix}0&1\\0&0\end{psmallmatrix}$。\\
Jordanian tensor $C_{12}$ & $C_{12}:=\Pup\otimes\Eplus-\Eplus\otimes\Pup$。\\
古典$r$-matrix $r_{12}(\zeta,\mu)$ & 空間Lax行列の線形Poisson algebraを与える$4\times4$行列。式\eqref{audit5r:eq:rcont}で定義する。\\
monodromy $T(x,y;\zeta)$ & $\partial_xT(x,y;\zeta)=U(x,\zeta)T(x,y;\zeta)$を満たすpath-ordered exponential。\\
局所monodromy projector $\Pi(x;\zeta)$ & $\zeta\to0$の局所漸近展開で、$\rho$へ近づくrank-one projector。式\eqref{audit5r:eq:Piexp}で定義する。\\
$r$-matrix routeの時間行列 $V_r$ & transfer matrixの一次のlocal chargeが生成する時間Lax行列。\\
量子極限の時間行列 $V_q$ & 前報告で有限格子の量子time operatorから直接得た連続時間Lax行列。\\
\bottomrule
\end{longtable}

\subsection{Class 5の連続空間Lax行列}

複素spin成分を
\begin{equation}
p=S^1+\ii S^2,\qquad q=S^1-\ii S^2,\qquad z=S^3,
\qquad pq+z^2=1
\label{audit5r:eq:spinconstraint}
\end{equation}
とする。spin-$\tfrac12$ coherent stateのrank-one projectorを
\begin{equation}
\rho(x)=\frac12
\begin{pmatrix}
1+z&q\\
p&1-z
\end{pmatrix},
\qquad
\rho^2=\rho,
\qquad
\tr\rho=1
\label{audit5r:eq:rho}
\end{equation}
と定義する。

前報告で量子$R$行列の長波長極限から得たClass 5の空間Lax行列は
\begin{equation}
U(x,\mu)
=
\frac{1}{4\mu}
\begin{pmatrix}
1+z+2\alpha\mu p&q-2\alpha\mu(1+z)\\
p&1-z
\end{pmatrix}
\label{audit5r:eq:Uexplicit}
\end{equation}
である。ここで連続変形パラメータを$\alpha$とした。

式\eqref{audit5r:eq:Uexplicit}を、小さい$\mu$で発散する部分と有限なJordanian変形に分けるため、有限項$K(x)$を
\begin{equation}
K(x)
:=\frac12
\begin{pmatrix}
p&-(1+z)\\0&0\end{pmatrix}
\label{audit5r:eq:K}
\end{equation}
と定義する。すると
\begin{equation}
U(x,\mu)=\frac{\rho(x)}{2\mu}+\alpha K(x)
\label{audit5r:eq:UK}
\end{equation}
と書ける。この分解が以下の小さい生成スペクトルパラメータ$\zeta$の展開を簡単にする。

Poisson bracketの規格化は前報告と同じく
\begin{equation}
\{p(x),q(y)\}=-4\ii z(x)\delta(x-y),\quad
\{p(x),z(y)\}=2\ii p(x)\delta(x-y),\quad
\{q(x),z(y)\}=-2\ii q(x)\delta(x-y)
\label{audit5r:eq:PBpqz}
\end{equation}
とする。

\subsection{連続古典$r$-matrixの固定}

Class 5原論文では、古典Lax行列がquadraticなKostant--Kirillov型Poisson bracketを満たし、その$r$-matrixを量子$R$行列から与えている\cite{DFN}。本報告では長波長規格化後の空間Lax行列$U$に直接Poisson bracketを取り、連続線形algebraの$r$-matrixを固定する。

まず、二つの補助空間を交換するpermutation作用素を
\begin{equation}
P_{12}=
\begin{pmatrix}
1&0&0&0\\
0&0&1&0\\
0&1&0&0\\
0&0&0&1
\end{pmatrix}
\end{equation}
とする。さらに、上向き基底へのprojectorとraising matrixを
\begin{equation}
\Pup=\begin{pmatrix}1&0\\0&0\end{pmatrix},
\qquad
\Eplus=\begin{pmatrix}0&1\\0&0\end{pmatrix}
\end{equation}
とし、Class 5のJordanian tensorを
\begin{equation}
C_{12}:=\Pup\otimes\Eplus-\Eplus\otimes\Pup
\label{audit5r:eq:Cj}
\end{equation}
と定義する。

このとき、前節のPoisson bracket規約に対応する連続古典$r$-matrixは
\begin{equation}
r_{12}(\zeta,\mu)
=
\frac{\ii}{2(\zeta-\mu)}P_{12}
+\ii\alpha C_{12}
\label{audit5r:eq:rcont}
\end{equation}
である。実際、補助空間1と2に作用する空間Lax行列を
\[
U_1=U\otimes\Id_2,
\qquad
U_2=\Id_2\otimes U
\]
と書くと、成分を直接計算して
\begin{equation}
\{U_1(x,\zeta),U_2(y,\mu)\}
=
[r_{12}(\zeta,\mu),U_1(x,\zeta)+U_2(y,\mu)]\,\delta(x-y)
\label{audit5r:eq:linearPB}
\end{equation}
が成立する。

式\eqref{audit5r:eq:rcont}の第一項は通常のisotropic Landau--Lifshitz模型のrational $r$-matrixに対応し、第二項がClass 5のJordanian変形である。原論文の式(2.36)にも、量子$R$行列のClass 5変形に由来する同じrational部分とJordanian部分が現れる\cite{DFN}。ここでは係数と$\ii$の位置を原論文から読み替えるのではなく、式\eqref{audit5r:eq:PBpqz}と式\eqref{audit5r:eq:Uexplicit}から直接固定した。

\subsection{monodromyから時間成分を生成する標準式}

空間Lax行列$U(x,\zeta)$から、区間$[y,x]$のmonodromyを
\begin{equation}
T(x,y;\zeta)
=
\mathcal P\exp\!\left(\int_y^x U(s,\zeta)\,ds\right)
\label{audit5r:eq:monodromy}
\end{equation}
と定義する。周期系ではtransfer matrixは
\begin{equation}
t(\zeta):=\tr T(L,0;\zeta)
\end{equation}
である。

線形Poisson algebra\eqref{audit5r:eq:linearPB}に対する標準的なmonodromy計算から、生成時間行列は
\begin{equation}
\mathbb V_b(x;\zeta,\mu)
=
t(\zeta)^{-1}
\tr_a\!\left[
T_a(L,x;\zeta)
\,r_{ab}(\zeta,\mu)\,
T_a(x,0;\zeta)
\right]
\label{audit5r:eq:VSTS}
\end{equation}
となる\cite{DK}。添字$a$はtraceを取る補助空間、添字$b$は最終的に$2\times2$行列として残る補助空間を表す。

ここで重要なのは、$\zeta$と$\mu$を区別することである。生成スペクトルパラメータ$\zeta$について$\ln t(\zeta)$と$\mathbb V(x;\zeta,\mu)$を展開し、その各係数が別々のHamiltonian flowを与える。一方、Laxスペクトルパラメータ$\mu$はそのflowのLax pairに残る。

\subsection{小さい生成スペクトルパラメータの局所projector}

Class 5の空間Lax行列は
\[
U(x,\zeta)=\frac{\rho(x)}{2\zeta}+\alpha K(x)
\]
なので、$\zeta\to0$ではrank-one projector $\rho$に比例する$1/\zeta$極が支配的である。monodromyのlocal chargeを得る標準的な漸近展開では、このleading eigenspaceに対応する局所rank-one projectorを用いる。

そのprojectorを$\Pi(x;\zeta)$と書き、
\begin{equation}
\Pi^2=\Pi,
\qquad
\tr\Pi=1,
\qquad
\partial_x\Pi=[U(x,\zeta),\Pi]
\label{audit5r:eq:Piconstraints}
\end{equation}
を満たし、$\zeta\to0$で$\rho$へ近づくbranchを選ぶ。展開を
\begin{equation}
\Pi(x;\zeta)=\rho(x)+\zeta\Pi_1(x)+O(\zeta^2)
\label{audit5r:eq:Piexp}
\end{equation}
と定義する。

式\eqref{audit5r:eq:Piconstraints}のprojector条件を$O(\zeta)$まで展開すると
\begin{equation}
\rho\Pi_1+\Pi_1\rho=\Pi_1
\label{audit5r:eq:Pi1projector}
\end{equation}
となる。一方、輸送方程式の$O(\zeta^0)$は
\begin{equation}
\rho_x
=
\frac12[\rho,\Pi_1]
+\alpha[K,\rho]
\label{audit5r:eq:Pi1transport}
\end{equation}
を与える。

$\rho$がrank-one projectorであることを使うと、これら二式の解は
\begin{equation}
\Pi_1
=
2[\rho,\rho_x]
-2\alpha[\rho,[K,\rho]]
\label{audit5r:eq:Pi1}
